%%%%%%%%%%%%%%%%% 图片浮动设置 %%%%%%%%%%%%%%%%%
\setcounter{topnumber}{2}
\setcounter{bottomnumber}{2}
\setcounter{totalnumber}{4}
\renewcommand{\topfraction}{0.85}
\renewcommand{\bottomfraction}{0.85}
\renewcommand{\textfraction}{0.15}
\renewcommand{\floatpagefraction}{0.8}
\renewcommand{\textfraction}{0.1}
\setlength{\floatsep}{5pt plus 2pt minus 2pt}
\setlength{\textfloatsep}{5pt plus 2pt minus 2pt}
\setlength{\intextsep}{5pt plus 2pt minus 2pt}

%%%%%%%%%%%%%%%%% Fonts Definition and Basics %%%%%%%%%%%%%%%%%
\newcommand{\song}{\CJKfamily{song}}    % 宋体
\newcommand{\fs}{\CJKfamily{fs}}        % 仿宋体
\newcommand{\kai}{\CJKfamily{kai}}      % 楷体
\newcommand{\hei}{\CJKfamily{hei}}      % 黑体
\newcommand{\li}{\CJKfamily{li}}        % 隶书
\newcommand{\chuhao}{\fontsize{28pt}{28pt}\selectfont}       % 初号, 单倍行距
\newcommand{\yihao}{\fontsize{26pt}{26pt}\selectfont}       % 一号, 单倍行距
\newcommand{\xiaoyi}{\fontsize{24pt}{24pt}\selectfont}      % 小一, 单倍行距
\newcommand{\erhao}{\fontsize{22pt}{1.25\baselineskip}\selectfont}       % 二号, 1.25倍行距
\newcommand{\xiaoer}{\fontsize{18pt}{18pt}\selectfont}      % 小二, 单倍行距
\newcommand{\sanhao}{\fontsize{16pt}{16pt}\selectfont}      % 三号, 单倍行距
\newcommand{\xiaosan}{\fontsize{15pt}{15pt}\selectfont}     % 小三, 单倍行距
\newcommand{\sihao}{\fontsize{14pt}{14pt}\selectfont}       % 四号, 单倍行距
\newcommand{\xiaosi}{\fontsize{12pt}{12pt}\selectfont}      % 小四, 单倍行距
\newcommand{\wuhao}{\fontsize{10.5pt}{10.5pt}\selectfont}   % 五号, 单倍行距
\newcommand{\xiaowu}{\fontsize{9pt}{9pt}\selectfont}        % 小五, 单倍行距

% 重新定义了波浪符~的意义
\CJKtilde

% 定义章的pre-post名称
\newcommand\prechaptername{第}
\newcommand\postchaptername{章}

% 调整中文字符的表示，行内占一个字符宽度，行尾占半个字符宽度
% 行末半角？
\punctstyle{hangmobanjiao}             

% 调整罗列环境的布局
\setitemize{leftmargin=3em,itemsep=0em,partopsep=0em,parsep=0em,topsep=-0em}
\setenumerate{leftmargin=3em,itemsep=0em,partopsep=0em,parsep=0em,topsep=0em,label={(\arabic*)}}

% 避免宏包 hyperref 和 arydshln 不兼容带来的目录链接失效的问题。
\def\temp{\relax}
\let\temp\addcontentsline
\gdef\addcontentsline{\phantomsection\temp}

% 自定义项目列表标签及格式 \begin{publist} 列表项 \end{publist}
\newcounter{pubctr} %自定义新计数器
\newenvironment{publist}{%%%%%定义新环境
\begin{list}{[\arabic{pubctr}]} %%标签格式
    {
     \usecounter{pubctr}
     \setlength{\leftmargin}{2.5em}   % 左边界 \leftmargin =\itemindent + \labelwidth + \labelsep
     \setlength{\itemindent}{0em}     % 标号缩进量
     \setlength{\labelsep}{1em}       % 标号和列表项之间的距离,默认0.5em
     \setlength{\rightmargin}{0em}    % 右边界
     \setlength{\topsep}{0ex}         % 列表到上下文的垂直距离
     \setlength{\parsep}{0ex}         % 段落间距
     \setlength{\itemsep}{0ex}        % 标签间距
     \setlength{\listparindent}{0pt}  % 段落缩进量
    }}
{\end{list}}

\makeatletter
	\renewcommand\normalsize{
		\@setfontsize\normalsize{12pt}{12pt} % 小四对应 12 pt
		\setlength\abovedisplayskip{4pt}
		\setlength\abovedisplayshortskip{4pt}
		\setlength\belowdisplayskip{\abovedisplayskip}
		\setlength\belowdisplayshortskip{\abovedisplayshortskip}
		\let\@listi\@listI}
	
	% 不同的行距设置
	% TJU原始值1.63
	% 设为1.8则一页31行，1.95则一页29行（目前采用值）（约1.5倍行距）
	\def\defaultfont{\renewcommand{\baselinestretch}{1.95}\normalsize\selectfont} % 设置行距，正文一页29行
	
	% 控制字间距，使每行 34 个汉字
	\renewcommand{\CJKglue}{\hskip -0.1 pt plus 0.08\baselineskip} 
\makeatother

%%%%%%%%%%%%% Contents 目录 %%%%%%%%%%%%%%%%%

\renewcommand{\contentsname}{目录}

% 控制目录深度，改为1
\setcounter{tocdepth}{1}

\titlecontents{chapter}[2em]{\vspace{.0\baselineskip}\sihao\song}	% 可以重调skip
	{\prechaptername\CJKnumber{\thecontentslabel}\postchaptername\quad}{}
	{\!\titlerule*[5pt]{$\cdot$}\!\!\!\!\sihao\contentspage}	% 调整点的距离
\titlecontents{section}[3em]{\vspace{-0.1\baselineskip}\xiaosi\song}
	{\thecontentslabel\quad}{}
	{\!\titlerule*[5pt]{$\cdot$}\!\!\!\!\xiaosi\contentspage}
\titlecontents{subsection}[4em]{\vspace{-0.2\baselineskip}\wuhao\song}
	{\thecontentslabel\quad}{}
	{\!\titlerule*[5pt]{$\cdot$}\!\!\!\!\wuhao\contentspage}
             
%%%%%%%%%% Chapter and Section 章节 %%%%%%%%%%%%%

\setcounter{secnumdepth}{4}
\setlength{\parindent}{2em}

% 如果使用第“一”章
\renewcommand{\chaptername}{\prechaptername\CJKnumber{\thechapter}\postchaptername}
% 使用第“1”章
%\renewcommand{\chaptername}{\thechapter}

% 此处修改的chapter title会被主文件定义覆盖
% chapter标题格式：小二，黑体，居中
\titleformat{\chapter}{\centering\xiaoer\hei}{\chaptername}{2em}{}
\titlespacing{\chapter}{0pt}{0.1\baselineskip}{0.8\baselineskip}

% section标题格式：四号，黑体，左对齐
\titleformat{\section}{\sihao\hei}{\thesection}{1em}{}
\titlespacing{\section}{0pt}{0.15\baselineskip}{0.25\baselineskip}

% subsection标题格式：小四，黑体，左对齐
\titleformat{\subsection}{\xiaosi\hei}{\thesubsection}{1em}{}
\titlespacing{\subsection}{0pt}{0.1\baselineskip}{0.3\baselineskip}

% subsubsection标题格式：小四，宋体加粗，左对齐
\titleformat{\subsubsection}{\xiaosi\song\bfseries}{\thesubsubsection}{1em}{}
\titlespacing{\subsubsection}{0pt}{0.05\baselineskip}{0.1\baselineskip}

%%%%%%%%%% Table, Figure and Equation 图/表/公式 %%%%%%%%%%%%%%%%%

\renewcommand{\tablename}{表}
\renewcommand{\figurename}{图}

% 使图编号为 7-1 的格式
%\renewcommand{\thefigure}{\arabic{chapter}-\arabic{figure}}
%按顺序编号
%\renewcommand{\thefigure}{\arabic{figure}}
% 使子图编号为 a) 的格式
%\renewcommand{\thesubfigure}{\alph{subfigure})}
% 使子图编号为 (a) 的格式
\renewcommand{\thesubfigure}{(\alph{subfigure})}

% 使子表编号为 (a) 的格式
\renewcommand{\thesubtable}{(\alph{subtable})}
% 使表编号为 7-1 的格式
%\renewcommand{\thetable}{\arabic{chapter}-\arabic{table}}
%按顺序编号
\renewcommand{\thetable}{\arabic{table}}
% 使公式编号为 7-1 的格式
%\renewcommand{\theequation}{\arabic{chapter}-\arabic{equation}}

\makeatletter
	% 使子图引用也是7-1a)或7-1(a)的形式
	\renewcommand{\p@subfigure}{\thefigure}
\makeatother

% 定制浮动图形和表格标题样式
\makeatletter
	\long\def\@makecaption#1#2{
	   \vskip\abovecaptionskip
	   \sbox\@tempboxa{\centering\wuhao\song\bfseries{#1\quad #2}}
	   \ifdim \wd\@tempboxa >\hsize
	     \centering\wuhao\song{#1\quad #2} \par	% narrower
	   \else
	     \global \@minipagefalse
	     \hb@xt@\hsize{\hfil\box\@tempboxa\hfil}
	   \fi
	   \vskip\belowcaptionskip}
\makeatother

% 用来控制longtable表头分隔符
\captiondelim{~~~~} 

%%%%%%%%%% Theorem Environment 定理 %%%%%%%%%%%%%%%%%
\theoremstyle{plain}
\theorembodyfont{\song\rmfamily}
\theoremheaderfont{\hei\rmfamily}
\newtheorem{theorem}{定理~}[chapter]
\newtheorem{lemma}{引理~}[chapter]
\newtheorem{axiom}{公理~}[chapter]
\newtheorem{proposition}{命题~}[chapter]
\newtheorem{prop}{性质~}[chapter]
\newtheorem{corollary}{推论~}[chapter]
\newtheorem{conclusion}{结论~}[chapter]
\newtheorem{definition}{定义~}[chapter]
\newtheorem{conjecture}{猜想~}[chapter]
\newtheorem{example}{例~}[chapter]
\newtheorem{remark}{注~}[chapter]
%\newtheorem{algorithm}{算法~}[chapter]
\newenvironment{proof}{\noindent{\hei 证明：}}{\hfill $ \square $ \vskip 4mm}
\theoremsymbol{$\square$}

%%%%%%%%%% Page: number, header and footer 页面设置 %%%%%%%%%%%%%%%%%

%\frontmatter 或 \pagenumbering{roman}
%\mainmatter 或 \pagenumbering{arabic}

\makeatletter
	\renewcommand\frontmatter{\clearpage
		\@mainmatterfalse}
\makeatother

%%%%%%%%%%%% References 参考文献 %%%%%%%%%%%%%%%%%

\renewcommand{\bibname}{参考文献}
% 重定义参考文献样式，来自thu
\makeatletter
\renewenvironment{thebibliography}[1]{
    %\titleformat{\chapter}{\raggedright\sihao\hei}{\chaptername}{2em}{}
    %\titleformat{\chapter}{\centering\sihao\hei}{\chaptername}{2em}{}
    %\titleformat{\chapter}{\centering\xiaoer\hei}{\chaptername}{2em}{}
   \chapter*{\bibname}
   \xiaosi
   \list{\@biblabel{\@arabic\c@enumiv}}
        {\renewcommand{\makelabel}[1]{##1\hfill}
         \settowidth\labelwidth{0 cm}
         \setlength{\labelsep}{0pt}
         \setlength{\itemindent}{0pt}
         \setlength{\leftmargin}{\labelwidth+\labelsep}
         \addtolength{\itemsep}{-0.7em}
%         \addtolength{\itemsep}{-1.0em}
         \linespread{1.5}\selectfont	% 调整每个参考文献项内的间距 ！！！
         \usecounter{enumiv}
         \let\p@enumiv\@empty
         \renewcommand\theenumiv{\@arabic\c@enumiv}}
    \sloppy\frenchspacing
    \clubpenalty4000
    \@clubpenalty \clubpenalty
    \widowpenalty4000
    \interlinepenalty4000
    \sfcode`\.\@m}
   {\def\@noitemerr
     {\@latex@warning{Empty `thebibliography' environment}}
    \endlist\frenchspacing}
\makeatother

% 缩小参考文献间的垂直间距
\addtolength{\bibsep}{-0.5em}

% 每个条目自第二行起缩进的距离
\setlength{\bibhang}{2em}

% 参考文献引用作为上标出现
\makeatletter
	\def\@cite#1#2{\textsuperscript{[{#1\if@tempswa , #2\fi}]}}
\makeatother

% 引用格式
\bibpunct{[}{]}{,}{s}{}{,}

%%%%%%%%%%%% Cover 封面、摘要、版权、致谢格式定义 %%%%%%%%%%%%%%%%%
 
\makeatletter % 一直到结尾

\def\ctitle#1{\def\@ctitle{#1}}\def\@ctitle{}
\def\etitle#1{\def\@etitle{#1}}\def\@etitle{}
\def\subtitle#1{\def\@subtitle{#1}}\def\@subtitle{}
\def\esubject#1{\def\@esubject{#1}}\def\@esubject{}
\def\cauthor#1{\def\@cauthor{#1}}\def\@cauthor{}
\def\eauthor#1{\def\@eauthor{#1}}\def\@eauthor{}
\def\csupervisor#1{\def\@csupervisor{#1}}\def\@csupervisor{}
\def\esupervisor#1{\def\@esupervisor{#1}}\def\@esupervisor{}
\def\cdate#1{\def\@cdate{#1}}\def\@cdate{}
\long\def\cabstract#1{\long\def\@cabstract{#1}}\long\def\@cabstract{}
\long\def\eabstract#1{\long\def\@eabstract{#1}}\long\def\@eabstract{}
\def\ckeywords#1{\def\@ckeywords{#1}}\def\@ckeywords{}
\def\ekeywords#1{\def\@ekeywords{#1}}\def\@ekeywords{}
\def\cheading#1{\def\@cheading{#1}}\def\@cheading{}

\pagestyle{fancy}
  \fancyhf{}
  \fancyhead[C]{\song\wuhao \@cheading}  % 页眉
%  \lhead{\song\wuhao \@cheading}  % 左页眉
%  \rhead{\prechaptername\CJKnumber{\thechapter}\postchaptername}    % 右页眉
%  \rhead{\prechaptername~\thechapter~\postchaptername}    % 右页眉
  \fancyfoot[C]{\song\xiaowu ~\thepage~}
\newlength{\@title@width}

% 定义封面
\def\makecover{
   \phantomsection
    \pdfbookmark[-1]{\@ctitle}{ctitle}

\begin{titlepage}
%\vspace*{31.5pt}
\begin{center}
  %\vspace*{21pt}
\setlength{\parindent}{2em}
  \begin{figure}[!h]
  \flushleft
  \hspace*{4pt}
  \includegraphics[width=3.18cm]{icon.png}
  \end{figure}

  \vspace*{21pt}
  \begin{figure}[!h]
  \centering
  \includegraphics[width=11.34cm]{name.png}
  \end{figure}
  \vspace*{42pt}

  \song\erhao{\textbf{第三十一届“冯如杯”学生学术科技作品竞赛项目论文}}

  \vspace*{110pt}
  \song\erhao{\@ctitle}

  \begin{flushright}
  \sanhao{\textrm{\@subtitle}}
  \end{flushright}

 
   \vspace*{42pt}
    \setlength{\@title@width}{5cm}

 % \vspace*{60pt}

 %  \vspace*{21pt}
\vfill
\sanhao{\textbf{\@cdate}}
\end{center}
\end{titlepage}


%%%%%%%%%%%%%%%%%%% Abstract and Keywords 摘要和关键词 %%%%%%%%%%%%%%%%%%%%%%%

%中文摘要格式
\clearpage
\markboth{摘要}{摘要}
\pdfbookmark[0]{摘要}{cabstract}

% 摘要不加到目录中
%\addcontentsline{toc}{chapter}{摘要}

% 开始罗马数字编号
\setcounter{page}{1}
\pagenumbering{roman}
\thispagestyle{plain}


% 中文摘要：小二，黑体加粗，居中
\begin{center}
% \xiaosan\hei\bfseries 摘要
\xiaosan\hei 摘要
\end{center}

%\vspace{\baselineskip} % 新增摘要后空行

% 插入中文摘要
\song\defaultfont
\@cabstract
% \vspace{\baselineskip}

%\hangafter=1\hangindent=52.3pt\noindent
\hangafter=1\hangindent=52.3pt
{\song\xiaosi 关键词：} \@ckeywords
%\thispagestyle{empty}

% 英文摘要格式
\clearpage
 \markboth{Abstract}{Abstract}
 \pdfbookmark[0]{Abstract}{eabstract}

 % 摘要不加到目录中
 %\addcontentsline{toc}{chapter}{ABSTRACT}

 \thispagestyle{plain}

% ABSTRACT三号居中
 \begin{center}
 % \sanhao{\bf{Abstract}}
 \sanhao{Abstract}
 \end{center}
% \vspace{\baselineskip}

% 插入英文摘要
\@eabstract
% \vspace{\baselineskip}

% \hangafter=1\hangindent=60pt\noindent
\hangafter=1\hangindent=60pt
{\textbf{Keywords:}} \@ekeywords
\thispagestyle{plain}

}
\makeatother
